\documentclass[11pt,a4paper]{article}

% -------------------------------------------------
% Paquetes
% -------------------------------------------------
\usepackage[utf8]{inputenc}
\usepackage[T1]{fontenc}
\usepackage[spanish]{babel}
\usepackage{geometry}
\usepackage{enumitem}
\usepackage{hyperref}
\usepackage{xcolor}
\usepackage{listings}
\usepackage{booktabs}
\usepackage{array}
\usepackage{tabularx}

\geometry{
  top=2.5cm,
  bottom=2.5cm,
  left=2.7cm,
  right=2.7cm
}

\hypersetup{
  colorlinks=true,
  linkcolor=black,
  urlcolor=blue
}

\setlist[itemize]{noitemsep, topsep=4pt}
\setlist[enumerate]{noitemsep, topsep=4pt}

\lstset{
  basicstyle=\ttfamily\small,
  frame=single,
  breaklines=true,
  columns=fullflexible
}

\lstset{
  literate=
  {á}{{\'a}}1 {é}{{\'e}}1 {í}{{\'i}}1 {ó}{{\'o}}1 {ú}{{\'u}}1
  {Á}{{\'A}}1 {É}{{\'E}}1 {Í}{{\'I}}1 {Ó}{{\'O}}1 {Ú}{{\'U}}1
  {ñ}{{\~n}}1 {Ñ}{{\~N}}1 {¿}{{?`}}1 {¡}{{!`}}1
}

% -------------------------------------------------
% Documento
% -------------------------------------------------
\title{
  \textbf{Protocolo de Trabajo con Git, GitHub y Linear}\\
  \large Taller de Integración II y Taller de Integración IV
}
\author{}
\date{Segundo semestre 2026}

\begin{document}

\maketitle

\section{Objetivo}

Esta guía fija la forma en que los equipos de Taller de Integración II (TI2) y
Taller de Integración IV (TI4) trabajarán en el proyecto.

El flujo busca dos cosas concretas: que el trabajo se pueda seguir desde Linear
hasta GitHub y que cada equipo conserve sus responsabilidades técnicas.

El proyecto usa tres herramientas:

\begin{itemize}
  \item \textbf{Linear}: planificación, asignación y seguimiento del trabajo.
  \item \textbf{Git}: control de versiones.
  \item \textbf{GitHub}: alojamiento del código, Pull Requests, revisión de
    código e integración continua.
\end{itemize}

La división de responsabilidades es directa:

\begin{center}
  \textbf{Linear gestiona el trabajo y GitHub gestiona el código.}
\end{center}

% -------------------------------------------------

\section{Organización de los equipos}

El proyecto estará compuesto por dos equipos.

\subsection{Taller de Integración II}

TI2 tendrá seis estudiantes y se encargará de:

\begin{itemize}
  \item desarrollo de la aplicación web;
  \item desarrollo del backend y API mediante arquitectura por capas;
  \item implementación de los casos de uso y reglas de negocio compartidas;
  \item diseño e implementación de la base de datos;
  \item pruebas asociadas a sus componentes;
  \item despliegue de sus propios servicios;
  \item documentación técnica correspondiente.
\end{itemize}

\subsection{Taller de Integración IV}

TI4 tendrá tres estudiantes y se encargará de:

\begin{itemize}
  \item desarrollo de la aplicación móvil;
  \item planificación general del proyecto;
  \item preparación y priorización inicial del trabajo;
  \item coordinación entre ambos equipos;
  \item control de calidad;
  \item revisión de código;
  \item seguimiento de dependencias;
  \item apoyo técnico e inducción al equipo TI2 cuando sea necesario.
\end{itemize}

TI4 debe ayudar a TI2 a resolver sus tareas cada vez con más autonomía. En una
revisión, TI4 señala el problema y explica el criterio. La persona autora hace
el cambio en su propia rama.

% -------------------------------------------------

\section{Repositorio GitHub}

Todo el producto vivirá en \textbf{un único repositorio GitHub}. El monorepo
contendrá las aplicaciones, el backend y los paquetes compartidos.

\begin{lstlisting}
Proyecto
|
+-- apps/
|   +-- web/                 Aplicación web de TI2
|   +-- mobile/              Aplicación móvil de TI4
|
+-- convex/                  Backend por capas y base de datos compartidos
|   +-- presentation/        Funciones públicas de Convex
|   +-- application/         Casos de uso
|   +-- domain/              Reglas de negocio
|   +-- infrastructure/      Persistencia y servicios externos
+-- packages/                Código compartido cuando corresponda
+-- .github/workflows/       Integración continua
+-- package.json
+-- bun.lock
\end{lstlisting}

Cada equipo desarrolla, prueba y despliega los componentes que tiene asignados
dentro del monorepo.

Un repositorio común no elimina la separación de responsabilidades. TI2
mantiene la aplicación web y el backend. TI4 mantiene la aplicación móvil.

La integración entre los equipos pasa por la API pública de la capa de
Presentación del backend. Ambas aplicaciones consumen sus tipos generados. No
se copian implementaciones ni se mantienen contratos externos.

\subsection{Arquitectura por capas}

Las dependencias deben avanzar en esta dirección:

\begin{lstlisting}
Presentación
      |
      v
Aplicación
      |
      v
Dominio
      ^
      |
Infraestructura
\end{lstlisting}

En el backend, las funciones públicas de Convex reciben la solicitud y llaman
a un caso de uso. Los casos de uso coordinan la operación. El Dominio contiene
las reglas. Infraestructura conecta con la persistencia, la autenticación y los
servicios externos.

Una branch o Pull Request que cruce varias capas debe explicar qué cambia en
cada una. React, Mobile y las funciones públicas no deben resolver necesidades
de negocio con consultas directas a la base de datos ni con funciones extensas.

% -------------------------------------------------

\section{Uso de Linear}

Todo el proyecto usará \textbf{un único workspace de Linear}.

Dentro de este workspace existirán dos Teams:

\begin{itemize}
  \item \textbf{TI2}: Web, Backend y Base de Datos.
  \item \textbf{TI4}: Gestión y Aplicación Móvil.
\end{itemize}

Los dos Teams forman parte del mismo proyecto.

\subsection{Triage}

TI4 hará el \textit{triage} inicial. Para cada issue debe:

\begin{itemize}
  \item crear o revisar los issues necesarios;
  \item asignarlos al Team correspondiente;
  \item definir su prioridad;
  \item verificar que tengan una descripción suficiente;
  \item establecer criterios de aceptación;
  \item identificar dependencias entre ambos equipos;
  \item asignarlos al Cycle correspondiente.
\end{itemize}

Cuando un issue esté listo en el backlog, el equipo responsable organiza su
ejecución.

\subsection{Linear como fuente de verdad}

Toda tarea técnica que requiera trabajo significativo debe tener un issue en
Linear.

WhatsApp, Discord y otros canales sirven para conversar, no para asignar trabajo
de manera formal. Si una conversación genera una tarea, alguien debe registrarla
en Linear.

\subsection{Sub-issues}

El proyecto no usa \textbf{sub-issues} de Linear.

Para una tarea con varios pasos pequeños, usa una lista en la descripción del
issue.

Por ejemplo:

\begin{lstlisting}
TI2-37 Implementar autenticación

- Crear endpoint de login
- Validar credenciales
- Generar token
- Agregar pruebas
- Documentar endpoint
\end{lstlisting}

Si una parte del trabajo necesita asignación, seguimiento o revisión propios,
conviértela en un issue independiente.

Así se evita gastar issues del plan en tareas que no necesitan seguimiento
separado.

% -------------------------------------------------

\section{Cycles y Sprints}

Cada \textit{Cycle} de Linear representa un Sprint. No crearemos un Cycle por
semana.

La correspondencia es esta:

\begin{lstlisting}
Cycle 1 -> Sprint 1
Cycle 2 -> Sprint 2
Cycle 3 -> Sprint 3
Cycle 4 -> Sprint Final
\end{lstlisting}

Las semanas de cada Sprint quedan para seguimiento, reuniones y control de
avance.

Al cerrar un Cycle, el equipo revisa:

\begin{itemize}
  \item trabajo completado;
  \item trabajo pendiente;
  \item issues bloqueados;
  \item deuda técnica generada;
  \item dependencias entre TI2 y TI4;
  \item trabajo que debe trasladarse al siguiente Sprint.
\end{itemize}

% -------------------------------------------------

\section{Estados de los Issues}

El workflow tendrá pocos estados:

\begin{center}
  \texttt{Backlog}
  $\rightarrow$
  \texttt{Todo}
  $\rightarrow$
  \texttt{In Progress}
  $\rightarrow$
  \texttt{In Review}
  $\rightarrow$
  \texttt{Done}
\end{center}

También puedes usar:

\begin{center}
  \texttt{Canceled}
\end{center}

Los estados significan lo siguiente:

\begin{description}
  \item[Backlog:] trabajo identificado pero todavía no comprometido.
  \item[Todo:] trabajo seleccionado para el Cycle actual.
  \item[In Progress:] existe trabajo activo sobre el issue.
  \item[In Review:] existe uno o más Pull Requests pendientes de revisión.
  \item[Done:] todo el trabajo requerido por el issue fue integrado.
  \item[Canceled:] el trabajo dejó de ser necesario.
\end{description}

Cuando sea posible, la integración Linear--GitHub cambiará estos estados de
forma automática.

% -------------------------------------------------

\section{Integración de Linear con GitHub}

El repositorio del monorepo debe conectarse al workspace de Linear.

Cada issue conserva el Team responsable, TI2 o TI4. El identificador distingue
el área aunque todo el código viva en el mismo repositorio.

Los identificadores de Linear relacionan el código con el trabajo.

Por ejemplo:

\begin{lstlisting}
TI2-37
TI4-18
\end{lstlisting}

Un Pull Request puede referenciar el issue desde la branch, el título o la
descripción.

La trazabilidad queda así:

\begin{center}
  \textbf{
    Issue Linear
    $\rightarrow$
    Branch
    $\rightarrow$
    Commits
    $\rightarrow$
    Pull Request
    $\rightarrow$
    Review
    $\rightarrow$
    Merge
  }
\end{center}

% -------------------------------------------------

\section{Estrategia de ramas}

El repositorio usará \textbf{feature branches pequeñas y de corta duración}.

La única rama permanente es:

\begin{lstlisting}
main
\end{lstlisting}

No habrá una rama permanente \texttt{develop}. Tampoco usaremos Git Flow con
ramas permanentes de desarrollo, release o hotfix.

\subsection{Creación de una rama}

Antes de empezar una tarea:

\begin{enumerate}
  \item confirma que existe un issue en Linear;
  \item actualiza la rama \texttt{main};
  \item crea una rama desde \texttt{main};
  \item trabaja solo en el objetivo del issue.
\end{enumerate}

Flujo conceptual:

\begin{lstlisting}
main
 |
 +-- rama TI2-37-login-web
 |
 +-- rama TI2-42-auth
 |
 +-- rama TI4-18-login-mobile
\end{lstlisting}

% -------------------------------------------------

\section{Nombres de ramas}

Usa, de preferencia, el nombre de branch que genera Linear.

El nombre debe conservar:

\begin{itemize}
  \item nombre o identificador del desarrollador;
  \item identificador del issue de Linear;
  \item una descripción breve.
\end{itemize}

Ejemplo:

\begin{lstlisting}
vrivera/ti2-37-login
\end{lstlisting}

Si Linear genera un nombre demasiado largo, puedes abreviarlo.

Por ejemplo, en lugar de:

\begin{lstlisting}
vrivera/ti2-37-implementar-endpoint-de-autenticacion-de-usuarios
\end{lstlisting}

usa:

\begin{lstlisting}
vrivera/ti2-37-login
\end{lstlisting}

Conserva siempre el identificador de Linear.

No uses:

\begin{lstlisting}
login
feature1
cambios
prueba
rama-nueva
\end{lstlisting}

% -------------------------------------------------

\section{Tamaño de las ramas}

Una rama debe resolver \textbf{una intención concreta}.

Mantén las ramas pequeñas para facilitar:

\begin{itemize}
  \item revisión;
  \item pruebas;
  \item resolución de conflictos;
  \item integración continua;
  \item detección de errores;
  \item reversión de cambios.
\end{itemize}

No fijaremos un límite artificial de líneas modificadas. El criterio es que otra
persona pueda entender y revisar el Pull Request sin encontrarse con varias
funcionalidades independientes.

% -------------------------------------------------

\section{Pull Requests}

Todo cambio que llegue a \texttt{main} debe pasar por un \textbf{Pull Request}
en el repositorio del monorepo.

Nadie puede hacer push directo a \texttt{main}.

El Pull Request debe indicar:

\begin{itemize}
  \item qué problema resuelve;
  \item qué cambios introduce;
  \item cómo puede probarse;
  \item issue de Linear relacionado;
  \item cualquier consideración relevante para el reviewer.
\end{itemize}

Ejemplo:

\begin{lstlisting}
TI2-37: Implementar login de usuarios
\end{lstlisting}

La descripción debe dar el contexto suficiente para revisar el cambio sin
reconstruirlo desde el código.

% -------------------------------------------------

\section{Relación entre Issues y Pull Requests}

Un issue y un Pull Request no tienen que mantener una relación 1:1.

Un issue representa una \textbf{unidad de trabajo}.

Un Pull Request representa una \textbf{unidad de integración de código}.

Un issue puede requerir varios Pull Requests.

Ejemplo:

\begin{lstlisting}
TI2-42 Autenticación

    |
    +-- PR 1: modelo de usuario
    |
    +-- PR 2: endpoint login
    |
    +-- PR 3: refresh token
\end{lstlisting}

El issue termina cuando todo el trabajo necesario está integrado.

% -------------------------------------------------

\section{Stacked Pull Requests}

Si no puedes dividir una tarea en cambios independientes, encadena Pull
Requests o usa \textit{stacked PRs}.

Ejemplo:

\begin{lstlisting}
main
 |
 +-- PR A
      |
      +-- PR B
           |
           +-- PR C
\end{lstlisting}

Esta técnica es una excepción. No es el flujo predeterminado.

Como límite práctico:

\begin{center}
  \textbf{usa como máximo 3 stacked PRs; el límite absoluto es 4.}
\end{center}

Si una cadena supera cuatro Pull Requests dependientes, revisa la planificación
y divide la tarea de nuevo.

% -------------------------------------------------

\section{Revisión de código}

Otra persona debe revisar cada Pull Request antes de integrarlo.

Una revisión puede terminar en:

\begin{itemize}
  \item aprobación;
  \item comentarios;
  \item solicitud de cambios.
\end{itemize}

La revisión debe comprobar:

\begin{itemize}
  \item cumplimiento del issue;
  \item corrección;
  \item claridad del código;
  \item mantenibilidad;
  \item errores potenciales;
  \item pruebas;
  \item seguridad;
  \item respeto de la dirección de dependencias entre capas;
  \item ubicación correcta de las reglas de negocio;
  \item compatibilidad con otros componentes;
  \item contrato de API cuando corresponda.
\end{itemize}

% -------------------------------------------------

\section{Progresión de las revisiones de TI2}

TI2 está en una etapa inicial de formación, así que la supervisión cambiará
durante el semestre.

\subsection*{Etapa inicial}

Durante la primera etapa, integrantes de TI4 revisarán los Pull Requests de TI2.

TI4 puede pedir cambios, explicar el problema y sugerir una solución. La persona
autora hace los cambios en su propia rama.

\subsection*{Etapa intermedia}

En la etapa intermedia, TI2 empezará a revisar código entre pares.

TI4 seguirá revisando directamente los cambios relacionados con:

\begin{itemize}
  \item API;
  \item arquitectura;
  \item base de datos;
  \item seguridad;
  \item integración con móvil;
  \item cambios con impacto significativo.
\end{itemize}

\subsection*{Etapa avanzada}

En la etapa avanzada, TI2 debería revisar e integrar sus cambios con autonomía.

TI4 se concentrará en:

\begin{itemize}
  \item coordinación;
  \item QA;
  \item arquitectura;
  \item integración;
  \item resolución de dependencias.
\end{itemize}

% -------------------------------------------------

\section{Política de Merge}

Usaremos \textbf{Squash and Merge}. Así, cada Pull Request deja un solo commit
en \texttt{main}, aunque la branch tenga commits intermedios.

Después del merge:
\begin{itemize}
  \item conserva la rama de desarrollo;
  \item actualiza el issue relacionado;
  \item incluye la documentación asociada.
\end{itemize}

Las ramas integradas forman parte del historial de trabajo y de la evidencia de
participación de cada integrante. Desactiva en GitHub la eliminación automática
de branches después del merge.
% -------------------------------------------------

\section{Protección de \texttt{main}}

Protege la rama \texttt{main} con las opciones disponibles en GitHub.

Como mínimo, exige:

\begin{itemize}
  \item no permitir push directo;
  \item exigir Pull Request;
  \item exigir al menos una aprobación;
  \item exigir que las validaciones automáticas requeridas finalicen
    correctamente;
  \item exigir resolución de conversaciones pendientes antes del merge.
\end{itemize}

Estas reglas aplican a TI2 y TI4. El equipo TI4 también debe pasar por Pull
Request.

% -------------------------------------------------

\section{Integración Continua}

GitHub Actions ejecutará las verificaciones de cada Pull Request.

Las rutas modificadas determinan qué verificaciones se ejecutan. Como mínimo:

\begin{itemize}
  \item instalar las dependencias desde el \texttt{package.json} y el único
    \texttt{bun.lock} de la raíz;
  \item ejecutar el build de la aplicación o paquete afectado;
  \item formatter;
  \item linter;
  \item pruebas unitarias del Dominio y de los casos de uso;
  \item pruebas de integración de Infraestructura y de Presentación cuando
    corresponda.
\end{itemize}

Si un cambio afecta el backend, Web, Mobile o un paquete compartido, la CI
ejecuta las verificaciones de todos los componentes afectados.

No integres un Pull Request con la CI fallando, salvo una excepción justificada.

% -------------------------------------------------

\section{Integración entre TI2 y TI4}

TI2 implementará y desplegará los servicios compartidos con \textbf{Convex} y
arquitectura por capas. Convex ofrece los puntos de entrada, la ejecución y la
persistencia que usan las aplicaciones. Los casos de uso y las reglas de
negocio no deben acoplarse a Convex. Better Auth funcionará como adaptador de
Infraestructura.

El monorepo permite que TI2 y TI4 modifiquen una funcionalidad en el mismo Pull
Request cuando necesiten coordinarse. También pueden separar los cambios si la
revisión o el despliegue lo justifican. En ambos casos, relacionen los cambios
en Linear.

Cuando un cambio de TI2 afecte una funcionalidad de TI4:

\begin{itemize}
  \item relaciona el cambio con un issue de Linear;
  \item identifica al equipo afectado;
  \item comunica los cambios incompatibles antes de integrarlos;
  \item actualiza la documentación del servicio;
  \item indica si afecta Presentación, Aplicación, Dominio o Infraestructura;
  \item actualiza autenticación, tipos o estructura de datos si corresponde.
\end{itemize}

% -------------------------------------------------

\section{Cambios en el modelo de datos}

\textbf{Convex} administra los datos como parte de Infraestructura. Versiona
en \texttt{convex/} el esquema, las funciones, las relaciones y las estructuras
de datos. Mantén las entidades y reglas del negocio en Dominio, no en el
esquema ni en la interfaz. Toda configuración necesaria para ejecutar el
sistema debe quedar documentada y versionada.

Si un cambio del modelo afecta una funcionalidad de TI4, registra esa
dependencia entre equipos en Linear.

% -------------------------------------------------

\section{Dependencias entre equipos}

Si una tarea depende del trabajo del otro equipo, registra la dependencia en
Linear.

Ejemplo:

\begin{lstlisting}
TI4-25 Pantalla de recuperación de clave

Bloqueada por:

TI2-61 Endpoint de recuperación de clave
\end{lstlisting}

TI4 hará seguimiento de estas dependencias durante la coordinación del proyecto.

% -------------------------------------------------

\section{Definición de Done}

Un issue no termina porque ``el código funciona en el computador del
desarrollador''.

Marca una tarea como \texttt{Done} cuando:

\begin{itemize}
  \item cumple sus criterios de aceptación;
  \item el código fue integrado mediante Pull Request;
  \item las revisiones fueron resueltas;
  \item CI finalizó correctamente;
  \item las pruebas necesarias fueron realizadas;
  \item la documentación fue actualizada cuando corresponde;
  \item las dependencias entre capas respetan la arquitectura definida;
  \item no deja dependencias conocidas sin registrar;
  \item el cambio está disponible en \texttt{main}.
\end{itemize}

% -------------------------------------------------

\section{Flujo de trabajo completo}

El flujo estándar es:

\begin{enumerate}
  \item TI4 hace el triage inicial del trabajo.
  \item Linear asigna el issue al Team y al Cycle correspondientes.
  \item El integrante toma la tarea desde Linear.
  \item Identifica las capas afectadas y el punto de entrada de la
    funcionalidad.
  \item Actualiza localmente \texttt{main}.
  \item Crea desde Linear una branch asociada al issue dentro del monorepo.
  \item Implementa el cambio.
  \item Hace commits en su branch.
  \item Publica la branch en GitHub.
  \item Abre un Pull Request.
  \item Relaciona el Pull Request con el issue en Linear.
  \item Ejecuta CI sobre el cambio.
  \item Otro integrante revisa el código.
  \item Si hay observaciones, el autor modifica su propia branch.
  \item Cuando CI y la revisión están aprobadas, hace Squash and Merge.
  \item Conserva la branch utilizada en el repositorio.
  \item Actualiza el estado del trabajo en Linear.
  \item Despliega el cambio cuando corresponda.
\end{enumerate}

El flujo queda así:

\begin{lstlisting}
Linear
  |
  v
Issue
  |
  v
Branch pequeña
  |
  v
Desarrollo
  |
  v
Pull Request
  |
  +----> CI
  |
  +----> Code Review
  |
  v
Squash Merge
  |
  v
main
  |
  v
Deploy

La branch se conserva.
\end{lstlisting}

% -------------------------------------------------

\section{Stack tecnológico}

Ambos equipos deben conocer el stack tecnológico, aunque cada uno mantenga
componentes distintos. El proyecto usa:

\begin{table}[h]
  \centering
  \begin{tabularx}{\textwidth}{>{\bfseries}l X}
    \toprule
    Componente & Tecnología seleccionada \\
    \midrule
    Modelo arquitectónico & Arquitectura por capas \\
    Frontend Web & React + Vite + TanStack Router \\
    Servidor / Backend & Convex \\
    Base de Datos & Convex \\
    Autenticación & Better Auth \\
    Aplicación Móvil & Expo + React Native \\
    Monorepo y dependencias & Bun workspaces \\
    Integración Continua & GitHub Actions \\
    Gestión de Proyecto & Linear \\
    Control de Versiones & Git + GitHub \\
    \bottomrule
  \end{tabularx}
\end{table}

TI2 despliega y mantiene los servicios compartidos, incluidas las capas de
Aplicación, Dominio e Infraestructura implementadas con Convex y la integración
de Better Auth. TI4 integra y mantiene la Presentación Mobile y su
infraestructura cliente dentro del mismo monorepo. Las bibliotecas auxiliares y
las configuraciones internas pueden evolucionar, pero cualquier cambio en los
contratos o las dependencias entre equipos debe coordinarse antes.

% -------------------------------------------------

\section{Principios generales}

Durante el proyecto rigen estos principios:

\begin{enumerate}
  \item Registra todo trabajo relevante para que pueda seguirse.
  \item Nadie hace push directo a \texttt{main}.
  \item Mantén las ramas pequeñas y de corta duración.
  \item Haz que otra persona revise el código.
  \item TI4 ayuda a TI2 a ganar autonomía sin reemplazar su trabajo.
  \item Usa Linear para representar el estado del trabajo.
  \item Usa GitHub para representar el estado del código.
  \item Mantén todos los componentes principales del producto en el monorepo.
  \item Separa Presentación, Aplicación, Dominio e Infraestructura.
  \item Mantén las reglas de negocio fuera de las aplicaciones cliente.
  \item Mantén explícitos y versionados la API compartida y los paquetes
    internos.
  \item Haz visibles los problemas cuanto antes.
\end{enumerate}

% -------------------------------------------------

\section{Resumen}

El proyecto queda organizado mediante estas relaciones:

\begin{center}
  \textbf{Linear = planificación y seguimiento}
\end{center}

\begin{center}
  \textbf{GitHub = código, revisión, CI e integración}
\end{center}

\begin{center}
  \textbf{Monorepo = código compartido, trazable y coordinado}
\end{center}

El resultado esperado es un proceso que ambos equipos puedan repetir y mantener
sin depender de una sola persona.

\end{document}
